% ============================================================
%  SECCIÓN 6 y 7 — HARDWARE: DISEÑO DEL DISPOSITIVO
%  Tesis: Equipo de Ensayos de Seguridad Eléctrica Portátil
%  Autor: Valentina Ferrini
% ============================================================
%
% INSTRUCCIONES DE USO EN OVERLEAF
% ---------------------------------
% 1. Asegurate de tener en el preámbulo de tu documento principal:
%      \usepackage{booktabs}      % tablas de calidad
%      \usepackage{multirow}      % celdas combinadas en tablas
%      \usepackage{longtable}     % tablas que ocupan más de una página
%      \usepackage{array}         % control de columnas
%      \usepackage{graphicx}      % figuras
%      \usepackage{hyperref}      % hipervínculos
%      \usepackage{float}         % control de posición de figuras [H]
%      \usepackage[table]{xcolor} % colores en tablas
%      \usepackage{geometry}      % márgenes (ya probablemente lo tenés)
%      \usepackage[spanish]{babel}
%
% 2. Copiá el contenido de este archivo y pegalo en la sección
%    correspondiente de tu documento principal, o bien usá:
%      % ============================================================
%  SECCIÓN 6 y 7 — HARDWARE: DISEÑO DEL DISPOSITIVO
%  Tesis: Equipo de Ensayos de Seguridad Eléctrica Portátil
%  Autor: Valentina Ferrini
% ============================================================
%
% INSTRUCCIONES DE USO EN OVERLEAF
% ---------------------------------
% 1. Asegurate de tener en el preámbulo de tu documento principal:
%      \usepackage{booktabs}      % tablas de calidad
%      \usepackage{multirow}      % celdas combinadas en tablas
%      \usepackage{longtable}     % tablas que ocupan más de una página
%      \usepackage{array}         % control de columnas
%      \usepackage{graphicx}      % figuras
%      \usepackage{hyperref}      % hipervínculos
%      \usepackage{float}         % control de posición de figuras [H]
%      \usepackage[table]{xcolor} % colores en tablas
%      \usepackage{geometry}      % márgenes (ya probablemente lo tenés)
%      \usepackage[spanish]{babel}
%
% 2. Copiá el contenido de este archivo y pegalo en la sección
%    correspondiente de tu documento principal, o bien usá:
%      % ============================================================
%  SECCIÓN 6 y 7 — HARDWARE: DISEÑO DEL DISPOSITIVO
%  Tesis: Equipo de Ensayos de Seguridad Eléctrica Portátil
%  Autor: Valentina Ferrini
% ============================================================
%
% INSTRUCCIONES DE USO EN OVERLEAF
% ---------------------------------
% 1. Asegurate de tener en el preámbulo de tu documento principal:
%      \usepackage{booktabs}      % tablas de calidad
%      \usepackage{multirow}      % celdas combinadas en tablas
%      \usepackage{longtable}     % tablas que ocupan más de una página
%      \usepackage{array}         % control de columnas
%      \usepackage{graphicx}      % figuras
%      \usepackage{hyperref}      % hipervínculos
%      \usepackage{float}         % control de posición de figuras [H]
%      \usepackage[table]{xcolor} % colores en tablas
%      \usepackage{geometry}      % márgenes (ya probablemente lo tenés)
%      \usepackage[spanish]{babel}
%
% 2. Copiá el contenido de este archivo y pegalo en la sección
%    correspondiente de tu documento principal, o bien usá:
%      % ============================================================
%  SECCIÓN 6 y 7 — HARDWARE: DISEÑO DEL DISPOSITIVO
%  Tesis: Equipo de Ensayos de Seguridad Eléctrica Portátil
%  Autor: Valentina Ferrini
% ============================================================
%
% INSTRUCCIONES DE USO EN OVERLEAF
% ---------------------------------
% 1. Asegurate de tener en el preámbulo de tu documento principal:
%      \usepackage{booktabs}      % tablas de calidad
%      \usepackage{multirow}      % celdas combinadas en tablas
%      \usepackage{longtable}     % tablas que ocupan más de una página
%      \usepackage{array}         % control de columnas
%      \usepackage{graphicx}      % figuras
%      \usepackage{hyperref}      % hipervínculos
%      \usepackage{float}         % control de posición de figuras [H]
%      \usepackage[table]{xcolor} % colores en tablas
%      \usepackage{geometry}      % márgenes (ya probablemente lo tenés)
%      \usepackage[spanish]{babel}
%
% 2. Copiá el contenido de este archivo y pegalo en la sección
%    correspondiente de tu documento principal, o bien usá:
%      \input{seccion_hardware}
%
% ============================================================

% ------------------------------------------------------------
\section{Requisitos del Sistema a Desarrollar — Hardware}
\label{sec:requisitos_hw}
% ------------------------------------------------------------

El subsistema de hardware constituye el núcleo físico del equipo de ensayos.
Su función es ejecutar, de forma controlada y repetible, los procedimientos
de medición definidos por la norma IEC~60335, garantizando la seguridad
del operario y la trazabilidad de los resultados.

El dispositivo debe cubrir los siguientes ensayos:

\begin{itemize}
    \item Medición de consumo de potencia (tensión, corriente, potencia activa y factor de potencia).
    \item Medición de corriente de fuga pasiva.
    \item Medición de corriente de fuga activa.
    \item Medición de resistencia de puesta a tierra.
    \item Ensayo de rigidez dieléctrica (tensión resistida).
\end{itemize}

En el marco del presente proyecto se desarrolla el prototipo funcional
centrado en el \textbf{ensayo de consumo de potencia}, siendo los ensayos
restantes contemplados en la arquitectura del sistema para su incorporación
en etapas futuras.

Las medidas de seguridad implementadas en el hardware incluyen:
aislamiento galvánico entre etapas de alta y baja tensión,
desconexión automática ante valores fuera de rango,
indicadores visuales de estado y protocolos de paro de emergencia.

% ------------------------------------------------------------
\section{Metodología y Desarrollo del Hardware}
\label{sec:metodologia_hw}
% ------------------------------------------------------------

\subsection{Diseño del dispositivo}
\label{subsec:diseno_dispositivo}

\subsubsection{Criterios de selección de componentes}
\label{subsubsec:criterios_seleccion}

La selección de los componentes del equipo se realizó mediante un
\textbf{análisis de mercado} en el que se evaluaron los siguientes criterios:

\begin{enumerate}
    \item \textbf{Relación calidad-precio:} se priorizaron componentes que
          ofrecen prestaciones adecuadas para los rangos de medición exigidos
          por la norma IEC~60335, sin incurrir en costos de equipamiento
          industrial de alta gama.

    \item \textbf{Soporte de la comunidad y documentación:} se seleccionaron
          plataformas y módulos con amplia documentación técnica, bibliotecas
          de software activas y comunidades de desarrollo consolidadas, lo que
          reduce el riesgo de obsolescencia y facilita la resolución de
          problemas durante el desarrollo.

    \item \textbf{Disponibilidad en el mercado local e importación:}
          se verificó la disponibilidad de cada componente tanto en
          distribuidores nacionales como en plataformas de importación
          (Mercado Libre, AliExpress, DigiKey), asegurando la reposición
          del stock ante fallas o escalado del proyecto.
\end{enumerate}

\subsubsection{Descripción de componentes seleccionados}
\label{subsubsec:componentes}

A continuación se describen los componentes principales del dispositivo,
agrupados por función.

% - - - - - - - - - - - - - - - - - - - - -
\paragraph{Unidad de procesamiento — ESP32}

El microcontrolador seleccionado es el \textbf{ESP32} (Espressif Systems),
en su versión de módulo de desarrollo de 38 pines.
Se trata de un sistema en chip (SoC) de doble núcleo Xtensa LX6
a 240~MHz, con 520~KB de SRAM, conectividad Wi-Fi 802.11~b/g/n
y Bluetooth 4.2/BLE, y soporte nativo para los protocolos
UART, SPI, I\textsuperscript{2}C y PWM.

Su elección se justifica por la amplia disponibilidad en el mercado
argentino, el extenso ecosistema de bibliotecas compatibles con
Arduino y MicroPython, y por su capacidad de comunicarse simultáneamente
con múltiples módulos periféricos a través de UART, requisito
indispensable para la integración con los módulos PZEM-004T.
Adicionalmente, la conectividad inalámbrica incorporada abre la posibilidad
de comunicación remota en versiones futuras del equipo, sin requerir
hardware adicional.

% - - - - - - - - - - - - - - - - - - - - -
\paragraph{Módulos de medición de potencia — PZEM-004T v3.0}

Para la adquisición de parámetros eléctricos se utilizan tres módulos
\textbf{PZEM-004T v3.0} (Peacefair).
Cada módulo mide en tiempo real:

\begin{itemize}
    \item Tensión de línea: rango 80--260~V\textsubscript{AC},
          resolución 0{,}1~V, precisión $\pm$0{,}5\%.
    \item Corriente: rango 0--100~A (con transformador de corriente externo
          de pinza incluido), resolución 0{,}001~A, precisión $\pm$0{,}5\%.
    \item Potencia activa: rango 0--23~kW, resolución 0{,}1~W.
    \item Energía acumulada: rango 0--9999~kWh, resolución 1~Wh.
    \item Frecuencia: rango 45--65~Hz, resolución 0{,}1~Hz.
    \item Factor de potencia: rango 0{,}00--1{,}00, resolución 0{,}01.
\end{itemize}

La comunicación con el ESP32 se realiza mediante el protocolo
\textbf{UART/Modbus RTU} a 9600~baudios.
Cada módulo posee una dirección Modbus configurable (0x01--0xF7),
lo que permite conectar los tres módulos en un bus RS-485 compartido
y direccionarlos de forma individual desde el microcontrolador.

La elección de tres unidades responde a la arquitectura de conmutación
del equipo: uno para el ensayo de consumo nominal,
uno para la medición de fuga pasiva
y uno para la medición de fuga activa,
permitiendo la transición entre ensayos sin necesidad de reconexión manual.

% - - - - - - - - - - - - - - - - - - - - -
\paragraph{Sistema de conmutación — Contactores y módulo de relés}

La conmutación entre ensayos se realiza mediante dos tipos de actuadores:

\textbf{Contactores monofásicos 12~A (×4):}
se emplean contactores de 12~A con bobina de 220~V~AC,
configuración de contactos 2~NA + 2~NC.
Su función es la conmutación de la tensión de red (220~V~AC)
hacia las distintas etapas del equipo (alimentación del DUT,
circuito de rigidez dieléctrica, circuito de fuga).
La acción mecánica de apertura garantiza el aislamiento galvánico
efectivo entre etapas, lo que es un requisito de seguridad
fundamental cuando se trabaja con tensiones de ensayo elevadas
(hasta 2500~V en el ensayo de rigidez dieléctrica).

\textbf{Módulo de relés USB de 8 canales (×1):}
se utiliza un módulo de relés de 8 canales controlado por USB,
con relés de 5~V~DC y capacidad de conmutación de hasta 10~A~AC.
Este módulo gestiona señales de control de baja potencia:
activación de indicadores, habilitación de etapas auxiliares
y señalización de estado al software.
La comunicación por USB permite el control directo desde la PC
sin necesidad de pasar por el ESP32, simplificando la arquitectura
de control para funciones no críticas en tiempo.

% - - - - - - - - - - - - - - - - - - - - -
\paragraph{Transformadores de ensayo}

El equipo incorpora dos transformadores de propósito específico:

\textbf{Transformador de rigidez dieléctrica — 220~V / 1000--2500~V:}
transformador elevador con salida seleccionable entre 1000~V
(ensayo Clase~I) y 2500~V (ensayo Clase~II),
con corriente de cortocircuito limitada a 30~mA.
Esta limitación es un requisito de seguridad estipulado en la
norma IEC~60335: la corriente máxima admisible durante el ensayo
de tensión resistida es de 50~mA, y la restricción a 30~mA en
la fuente asegura que una falla del aislamiento no derive en
una corriente destructiva para el DUT ni peligrosa para el operario.

\textbf{Transformador de puesta a tierra — 220~V / 2{,}2~V / 75~VA:}
transformador reductor de 75~VA utilizado en el ensayo de
continuidad de tierra.
Entrega una tensión de 2{,}2~V~AC con capacidad de corriente
de ensayo de hasta 15~A~AC (equivalente a $75~\text{VA} / 5~\text{V} \approx 15~\text{A}$),
cumpliendo el requisito de la norma de inyectar al menos 25~A
o el 200\% de la corriente nominal del fusible de protección,
según el método aplicado.
La baja tensión de salida garantiza la seguridad del operario
durante el ensayo.

% - - - - - - - - - - - - - - - - - - - - -
\paragraph{Carcasa y estructura — Anvil Flight Case}

El equipo se aloja en un \textbf{Anvil Flight Case} de aluminio y
madera contrachapada con herrajes de acero.
Este tipo de contenedor, habitualmente utilizado en equipamiento
de audio y escenografía profesional, presenta las siguientes ventajas
para la aplicación:

\begin{itemize}
    \item Rigidez estructural adecuada para soportar transformadores
          y contactores sin deformaciones.
    \item Cierre con pestillos de seguridad que evitan apertura accidental.
    \item Carácter portátil con manija y esquinas reforzadas,
          compatible con el concepto de equipo de laboratorio portátil.
    \item Posibilidad de montaje de paneles frontales y rieles DIN
          en el interior para organizar los componentes.
\end{itemize}

% - - - - - - - - - - - - - - - - - - - - -
\paragraph{Placas de interfaz y conexionado}

\noindent\textbf{Nota para revisión:} \textit{Se incorporarán dos placas de
circuito impreso (PCB) de diseño propio cuyas especificaciones
se encuentran en proceso de definición (a confirmar con Nico).
Se completará esta sección con: función de cada placa,
dimensiones, descripción del circuito y justificación de diseño.}

% ------------------------------------------------------------
\subsection{Arquitectura del dispositivo}
\label{subsec:arquitectura_hw}
% ------------------------------------------------------------

La arquitectura del dispositivo se organiza en cinco bloques funcionales
interconectados a través de un bus de comunicación principal
gestionado por el ESP32:

\begin{enumerate}
    \item \textbf{Bloque de potencia:} transformadores, contactores y
          protecciones. Maneja las tensiones de red y de ensayo.
    \item \textbf{Bloque de adquisición:} módulos PZEM-004T. Realizan
          la medición de los parámetros eléctricos del DUT.
    \item \textbf{Bloque de control:} ESP32 y módulo de relés USB.
          Orchestran la secuencia de ensayos y la comunicación.
    \item \textbf{Bloque de comunicación:} interfaz USB/UART hacia la PC.
          Transmite datos crudos y recibe comandos de la interfaz gráfica.
    \item \textbf{Bloque de seguridad:} fusibles, protecciones térmicas
          e indicadores. Actúan de forma independiente ante condiciones
          de falla.
\end{enumerate}

% (Figura del diagrama de bloques — reemplazar con tu figura real)
\begin{figure}[H]
    \centering
    % \includegraphics[width=0.85\textwidth]{figuras/diagrama_bloques_hw.png}
    \fbox{\parbox{0.85\textwidth}{\centering\vspace{2cm}
    \textit{[Insertar aquí el diagrama de bloques del hardware]}
    \vspace{2cm}}}
    \caption{Diagrama de bloques del dispositivo de ensayos.}
    \label{fig:diagrama_bloques_hw}
\end{figure}

% ------------------------------------------------------------
\subsection{Protocolo de comunicación}
\label{subsec:protocolo_comunicacion}
% ------------------------------------------------------------

La comunicación entre componentes del equipo sigue el esquema
descrito a continuación:

\begin{itemize}
    \item \textbf{ESP32 $\leftrightarrow$ módulos PZEM-004T:}
          protocolo Modbus~RTU sobre UART a 9600~baudios, 8N1
          (8 bits de datos, sin paridad, 1 bit de parada).
          Los tres módulos comparten el bus serie;
          cada uno responde únicamente a su dirección Modbus asignada
          (0x01, 0x02 y 0x03 respectivamente).

    \item \textbf{ESP32 $\leftrightarrow$ PC:}
          puerto UART expuesto como dispositivo CDC-ACM (serial virtual)
          a través del chip USB-to-UART integrado en el módulo de desarrollo
          (CH340 o CP2102 según la revisión del módulo).
          La velocidad de comunicación es de 115200~baudios.
          El protocolo de capa de aplicación es un formato JSON
          delimitado por salto de línea (\texttt{newline-delimited JSON}),
          elegido por su legibilidad, facilidad de parseo en Python
          y extensibilidad para agregar nuevos tipos de medición.

    \item \textbf{PC $\leftrightarrow$ módulo de relés USB:}
          protocolo propietario del módulo sobre USB HID / CDC,
          controlado directamente desde el software en Python
          mediante la biblioteca \texttt{pyserial} o el driver
          correspondiente al chip del módulo.
\end{itemize}

% ------------------------------------------------------------
\section{Evaluación de costos}
\label{sec:costos}
% ------------------------------------------------------------

La siguiente tabla presenta el desglose de costos del prototipo
desarrollado. Los precios expresados en dólares estadounidenses (USD)
corresponden a valores de referencia de mercado obtenidos en
distribuidores nacionales y plataformas de importación
(Mercado Libre Argentina, AliExpress, DigiKey),
relevados durante el segundo semestre de 2025.
Los valores deben considerarse estimativos;
variaciones de tipo de cambio o disponibilidad de stock
pueden afectar el costo final.

\begin{table}[H]
\centering
\caption{Desglose de costos del prototipo — valores de referencia en USD}
\label{tab:costos_detalle}
\renewcommand{\arraystretch}{1.3}
\begin{tabular}{@{} l l r r r @{}}
\toprule
\textbf{Categoría} & \textbf{Componente} & \textbf{Cant.} & \textbf{P. Unit. (USD)} & \textbf{Subtotal (USD)} \\
\midrule
\multirow{2}{*}{Procesamiento}
  & ESP32 módulo de desarrollo (38 pines) & 1 & 8{,}00  & 8{,}00  \\
  & Módulo de relés USB 8 canales 5V DC   & 1 & 12{,}00 & 12{,}00 \\
\midrule
\multirow{3}{*}{Medición}
  & PZEM-004T v3.0 (con TC de pinza)      & 3 & 12{,}00 & 36{,}00 \\
  & \textit{Placas PCB de interfaz (×2)}  & 2 & —       & \textit{a confirmar} \\
  & Cables de medición y puntas           & — & 10{,}00 & 10{,}00 \\
\midrule
\multirow{2}{*}{Potencia}
  & Contactor 12A / 220V AC (2NA+2NC)     & 4 & 9{,}00  & 36{,}00 \\
  & Transformador 220V / 1000--2500V 30mA & 1 & 80{,}00 & 80{,}00 \\
  & Transformador 220V / 2{,}2V / 75VA    & 1 & 35{,}00 & 35{,}00 \\
\midrule
\multirow{3}{*}{Seguridad}
  & Fusible cerámico 10A + portafusible    & 4 & 2{,}00  & 8{,}00  \\
  & Piloto indicador LED 220V (rojo/verde) & 4 & 1{,}50  & 6{,}00  \\
  & Pulsador de paro de emergencia (hongo) & 1 & 8{,}00  & 8{,}00  \\
\midrule
\multirow{2}{*}{Conexionado}
  & Cable H07V-K 1,5~mm² (carrete 10~m)   & 2 & 6{,}00  & 12{,}00 \\
  & Borneras, terminales y bridas          & — & 8{,}00  & 8{,}00  \\
  & Toma corriente industrial + fichas     & 2 & 5{,}00  & 10{,}00 \\
\midrule
Estructura
  & Anvil Flight Case (aluminio)           & 1 & 60{,}00 & 60{,}00 \\
  & Riel DIN + panel frontal aluminio      & 1 & 15{,}00 & 15{,}00 \\
\midrule
\multicolumn{4}{r}{\textbf{Total estimado}} & \textbf{344{,}00} \\
\multicolumn{4}{r}{Margen de contingencia (15\%)} & 51{,}60 \\
\multicolumn{4}{r}{\textbf{Total con contingencia}} & \textbf{395{,}60} \\
\bottomrule
\end{tabular}
\end{table}

\noindent\textit{Nota: el costo de horas de ingeniería y desarrollo de
software no se incluye en el presupuesto del prototipo, dado que
corresponde al esfuerzo académico del proyecto de grado.}

% ------------------------------------------------------------
\subsection{Comparación con soluciones comerciales}
\label{subsec:comparacion_comercial}
% ------------------------------------------------------------

A modo de referencia, la siguiente tabla compara el equipo desarrollado
con equipos comerciales equivalentes disponibles en el mercado nacional,
en particular los modelos de la línea SYGlab (SYG~Group, Rosario, Argentina),
que realizan ensayos similares bajo norma IEC~60335.

\begin{table}[H]
\centering
\caption{Comparación entre el prototipo desarrollado y equipos comerciales equivalentes}
\label{tab:comparacion_comercial}
\renewcommand{\arraystretch}{1.3}
\begin{tabular}{@{} l c c c @{}}
\toprule
\textbf{Característica}       & \textbf{Prototipo (este trabajo)} & \textbf{SYGlab CL12-S0} & \textbf{SYGlab MG5} \\
\midrule
Ensayo de consumo / potencia  & \checkmark & \checkmark & \checkmark \\
Corriente de fuga pasiva      & Arq. prevista & \checkmark & \checkmark \\
Corriente de fuga activa      & Arq. prevista & \checkmark & \checkmark \\
Puesta a tierra               & Arq. prevista & \checkmark & \checkmark \\
Rigidez dieléctrica           & Arq. prevista & \checkmark & \checkmark \\
Clase I y II                  & \checkmark & \checkmark & \checkmark \\
Canales simultáneos           & 1          & 1          & 5          \\
Pantalla / GUI                & PC (PyQt)  & Pantalla HD integrada & Pantalla HD integrada \\
Exportación de datos          & CSV / JSON & Etiqueta QR + base de datos & Etiqueta QR + base de datos \\
Norma declarada               & IEC~60335  & IEC / IRAM & IEC / IRAM \\
Costo estimado                & $\sim$\$400 USD & Consultar  & Consultar  \\
\bottomrule
\end{tabular}
\end{table}

La comparación pone en evidencia que el prototipo desarrollado cubre
el mismo conjunto de ensayos que los equipos comerciales de referencia,
con la diferencia de que en este trabajo se implementa íntegramente
solo el ensayo de consumo, reservando los restantes para etapas
posteriores del proyecto.
La ventaja diferencial del prototipo radica en su bajo costo de
construcción, la apertura total del código fuente y la posibilidad
de adaptación a distintas normativas sin restricciones de licencia.

%
% ============================================================

% ------------------------------------------------------------
\section{Requisitos del Sistema a Desarrollar — Hardware}
\label{sec:requisitos_hw}
% ------------------------------------------------------------

El subsistema de hardware constituye el núcleo físico del equipo de ensayos.
Su función es ejecutar, de forma controlada y repetible, los procedimientos
de medición definidos por la norma IEC~60335, garantizando la seguridad
del operario y la trazabilidad de los resultados.

El dispositivo debe cubrir los siguientes ensayos:

\begin{itemize}
    \item Medición de consumo de potencia (tensión, corriente, potencia activa y factor de potencia).
    \item Medición de corriente de fuga pasiva.
    \item Medición de corriente de fuga activa.
    \item Medición de resistencia de puesta a tierra.
    \item Ensayo de rigidez dieléctrica (tensión resistida).
\end{itemize}

En el marco del presente proyecto se desarrolla el prototipo funcional
centrado en el \textbf{ensayo de consumo de potencia}, siendo los ensayos
restantes contemplados en la arquitectura del sistema para su incorporación
en etapas futuras.

Las medidas de seguridad implementadas en el hardware incluyen:
aislamiento galvánico entre etapas de alta y baja tensión,
desconexión automática ante valores fuera de rango,
indicadores visuales de estado y protocolos de paro de emergencia.

% ------------------------------------------------------------
\section{Metodología y Desarrollo del Hardware}
\label{sec:metodologia_hw}
% ------------------------------------------------------------

\subsection{Diseño del dispositivo}
\label{subsec:diseno_dispositivo}

\subsubsection{Criterios de selección de componentes}
\label{subsubsec:criterios_seleccion}

La selección de los componentes del equipo se realizó mediante un
\textbf{análisis de mercado} en el que se evaluaron los siguientes criterios:

\begin{enumerate}
    \item \textbf{Relación calidad-precio:} se priorizaron componentes que
          ofrecen prestaciones adecuadas para los rangos de medición exigidos
          por la norma IEC~60335, sin incurrir en costos de equipamiento
          industrial de alta gama.

    \item \textbf{Soporte de la comunidad y documentación:} se seleccionaron
          plataformas y módulos con amplia documentación técnica, bibliotecas
          de software activas y comunidades de desarrollo consolidadas, lo que
          reduce el riesgo de obsolescencia y facilita la resolución de
          problemas durante el desarrollo.

    \item \textbf{Disponibilidad en el mercado local e importación:}
          se verificó la disponibilidad de cada componente tanto en
          distribuidores nacionales como en plataformas de importación
          (Mercado Libre, AliExpress, DigiKey), asegurando la reposición
          del stock ante fallas o escalado del proyecto.
\end{enumerate}

\subsubsection{Descripción de componentes seleccionados}
\label{subsubsec:componentes}

A continuación se describen los componentes principales del dispositivo,
agrupados por función.

% - - - - - - - - - - - - - - - - - - - - -
\paragraph{Unidad de procesamiento — ESP32}

El microcontrolador seleccionado es el \textbf{ESP32} (Espressif Systems),
en su versión de módulo de desarrollo de 38 pines.
Se trata de un sistema en chip (SoC) de doble núcleo Xtensa LX6
a 240~MHz, con 520~KB de SRAM, conectividad Wi-Fi 802.11~b/g/n
y Bluetooth 4.2/BLE, y soporte nativo para los protocolos
UART, SPI, I\textsuperscript{2}C y PWM.

Su elección se justifica por la amplia disponibilidad en el mercado
argentino, el extenso ecosistema de bibliotecas compatibles con
Arduino y MicroPython, y por su capacidad de comunicarse simultáneamente
con múltiples módulos periféricos a través de UART, requisito
indispensable para la integración con los módulos PZEM-004T.
Adicionalmente, la conectividad inalámbrica incorporada abre la posibilidad
de comunicación remota en versiones futuras del equipo, sin requerir
hardware adicional.

% - - - - - - - - - - - - - - - - - - - - -
\paragraph{Módulos de medición de potencia — PZEM-004T v3.0}

Para la adquisición de parámetros eléctricos se utilizan tres módulos
\textbf{PZEM-004T v3.0} (Peacefair).
Cada módulo mide en tiempo real:

\begin{itemize}
    \item Tensión de línea: rango 80--260~V\textsubscript{AC},
          resolución 0{,}1~V, precisión $\pm$0{,}5\%.
    \item Corriente: rango 0--100~A (con transformador de corriente externo
          de pinza incluido), resolución 0{,}001~A, precisión $\pm$0{,}5\%.
    \item Potencia activa: rango 0--23~kW, resolución 0{,}1~W.
    \item Energía acumulada: rango 0--9999~kWh, resolución 1~Wh.
    \item Frecuencia: rango 45--65~Hz, resolución 0{,}1~Hz.
    \item Factor de potencia: rango 0{,}00--1{,}00, resolución 0{,}01.
\end{itemize}

La comunicación con el ESP32 se realiza mediante el protocolo
\textbf{UART/Modbus RTU} a 9600~baudios.
Cada módulo posee una dirección Modbus configurable (0x01--0xF7),
lo que permite conectar los tres módulos en un bus RS-485 compartido
y direccionarlos de forma individual desde el microcontrolador.

La elección de tres unidades responde a la arquitectura de conmutación
del equipo: uno para el ensayo de consumo nominal,
uno para la medición de fuga pasiva
y uno para la medición de fuga activa,
permitiendo la transición entre ensayos sin necesidad de reconexión manual.

% - - - - - - - - - - - - - - - - - - - - -
\paragraph{Sistema de conmutación — Contactores y módulo de relés}

La conmutación entre ensayos se realiza mediante dos tipos de actuadores:

\textbf{Contactores monofásicos 12~A (×4):}
se emplean contactores de 12~A con bobina de 220~V~AC,
configuración de contactos 2~NA + 2~NC.
Su función es la conmutación de la tensión de red (220~V~AC)
hacia las distintas etapas del equipo (alimentación del DUT,
circuito de rigidez dieléctrica, circuito de fuga).
La acción mecánica de apertura garantiza el aislamiento galvánico
efectivo entre etapas, lo que es un requisito de seguridad
fundamental cuando se trabaja con tensiones de ensayo elevadas
(hasta 2500~V en el ensayo de rigidez dieléctrica).

\textbf{Módulo de relés USB de 8 canales (×1):}
se utiliza un módulo de relés de 8 canales controlado por USB,
con relés de 5~V~DC y capacidad de conmutación de hasta 10~A~AC.
Este módulo gestiona señales de control de baja potencia:
activación de indicadores, habilitación de etapas auxiliares
y señalización de estado al software.
La comunicación por USB permite el control directo desde la PC
sin necesidad de pasar por el ESP32, simplificando la arquitectura
de control para funciones no críticas en tiempo.

% - - - - - - - - - - - - - - - - - - - - -
\paragraph{Transformadores de ensayo}

El equipo incorpora dos transformadores de propósito específico:

\textbf{Transformador de rigidez dieléctrica — 220~V / 1000--2500~V:}
transformador elevador con salida seleccionable entre 1000~V
(ensayo Clase~I) y 2500~V (ensayo Clase~II),
con corriente de cortocircuito limitada a 30~mA.
Esta limitación es un requisito de seguridad estipulado en la
norma IEC~60335: la corriente máxima admisible durante el ensayo
de tensión resistida es de 50~mA, y la restricción a 30~mA en
la fuente asegura que una falla del aislamiento no derive en
una corriente destructiva para el DUT ni peligrosa para el operario.

\textbf{Transformador de puesta a tierra — 220~V / 2{,}2~V / 75~VA:}
transformador reductor de 75~VA utilizado en el ensayo de
continuidad de tierra.
Entrega una tensión de 2{,}2~V~AC con capacidad de corriente
de ensayo de hasta 15~A~AC (equivalente a $75~\text{VA} / 5~\text{V} \approx 15~\text{A}$),
cumpliendo el requisito de la norma de inyectar al menos 25~A
o el 200\% de la corriente nominal del fusible de protección,
según el método aplicado.
La baja tensión de salida garantiza la seguridad del operario
durante el ensayo.

% - - - - - - - - - - - - - - - - - - - - -
\paragraph{Carcasa y estructura — Anvil Flight Case}

El equipo se aloja en un \textbf{Anvil Flight Case} de aluminio y
madera contrachapada con herrajes de acero.
Este tipo de contenedor, habitualmente utilizado en equipamiento
de audio y escenografía profesional, presenta las siguientes ventajas
para la aplicación:

\begin{itemize}
    \item Rigidez estructural adecuada para soportar transformadores
          y contactores sin deformaciones.
    \item Cierre con pestillos de seguridad que evitan apertura accidental.
    \item Carácter portátil con manija y esquinas reforzadas,
          compatible con el concepto de equipo de laboratorio portátil.
    \item Posibilidad de montaje de paneles frontales y rieles DIN
          en el interior para organizar los componentes.
\end{itemize}

% - - - - - - - - - - - - - - - - - - - - -
\paragraph{Placas de interfaz y conexionado}

\noindent\textbf{Nota para revisión:} \textit{Se incorporarán dos placas de
circuito impreso (PCB) de diseño propio cuyas especificaciones
se encuentran en proceso de definición (a confirmar con Nico).
Se completará esta sección con: función de cada placa,
dimensiones, descripción del circuito y justificación de diseño.}

% ------------------------------------------------------------
\subsection{Arquitectura del dispositivo}
\label{subsec:arquitectura_hw}
% ------------------------------------------------------------

La arquitectura del dispositivo se organiza en cinco bloques funcionales
interconectados a través de un bus de comunicación principal
gestionado por el ESP32:

\begin{enumerate}
    \item \textbf{Bloque de potencia:} transformadores, contactores y
          protecciones. Maneja las tensiones de red y de ensayo.
    \item \textbf{Bloque de adquisición:} módulos PZEM-004T. Realizan
          la medición de los parámetros eléctricos del DUT.
    \item \textbf{Bloque de control:} ESP32 y módulo de relés USB.
          Orchestran la secuencia de ensayos y la comunicación.
    \item \textbf{Bloque de comunicación:} interfaz USB/UART hacia la PC.
          Transmite datos crudos y recibe comandos de la interfaz gráfica.
    \item \textbf{Bloque de seguridad:} fusibles, protecciones térmicas
          e indicadores. Actúan de forma independiente ante condiciones
          de falla.
\end{enumerate}

% (Figura del diagrama de bloques — reemplazar con tu figura real)
\begin{figure}[H]
    \centering
    % \includegraphics[width=0.85\textwidth]{figuras/diagrama_bloques_hw.png}
    \fbox{\parbox{0.85\textwidth}{\centering\vspace{2cm}
    \textit{[Insertar aquí el diagrama de bloques del hardware]}
    \vspace{2cm}}}
    \caption{Diagrama de bloques del dispositivo de ensayos.}
    \label{fig:diagrama_bloques_hw}
\end{figure}

% ------------------------------------------------------------
\subsection{Protocolo de comunicación}
\label{subsec:protocolo_comunicacion}
% ------------------------------------------------------------

La comunicación entre componentes del equipo sigue el esquema
descrito a continuación:

\begin{itemize}
    \item \textbf{ESP32 $\leftrightarrow$ módulos PZEM-004T:}
          protocolo Modbus~RTU sobre UART a 9600~baudios, 8N1
          (8 bits de datos, sin paridad, 1 bit de parada).
          Los tres módulos comparten el bus serie;
          cada uno responde únicamente a su dirección Modbus asignada
          (0x01, 0x02 y 0x03 respectivamente).

    \item \textbf{ESP32 $\leftrightarrow$ PC:}
          puerto UART expuesto como dispositivo CDC-ACM (serial virtual)
          a través del chip USB-to-UART integrado en el módulo de desarrollo
          (CH340 o CP2102 según la revisión del módulo).
          La velocidad de comunicación es de 115200~baudios.
          El protocolo de capa de aplicación es un formato JSON
          delimitado por salto de línea (\texttt{newline-delimited JSON}),
          elegido por su legibilidad, facilidad de parseo en Python
          y extensibilidad para agregar nuevos tipos de medición.

    \item \textbf{PC $\leftrightarrow$ módulo de relés USB:}
          protocolo propietario del módulo sobre USB HID / CDC,
          controlado directamente desde el software en Python
          mediante la biblioteca \texttt{pyserial} o el driver
          correspondiente al chip del módulo.
\end{itemize}

% ------------------------------------------------------------
\section{Evaluación de costos}
\label{sec:costos}
% ------------------------------------------------------------

La siguiente tabla presenta el desglose de costos del prototipo
desarrollado. Los precios expresados en dólares estadounidenses (USD)
corresponden a valores de referencia de mercado obtenidos en
distribuidores nacionales y plataformas de importación
(Mercado Libre Argentina, AliExpress, DigiKey),
relevados durante el segundo semestre de 2025.
Los valores deben considerarse estimativos;
variaciones de tipo de cambio o disponibilidad de stock
pueden afectar el costo final.

\begin{table}[H]
\centering
\caption{Desglose de costos del prototipo — valores de referencia en USD}
\label{tab:costos_detalle}
\renewcommand{\arraystretch}{1.3}
\begin{tabular}{@{} l l r r r @{}}
\toprule
\textbf{Categoría} & \textbf{Componente} & \textbf{Cant.} & \textbf{P. Unit. (USD)} & \textbf{Subtotal (USD)} \\
\midrule
\multirow{2}{*}{Procesamiento}
  & ESP32 módulo de desarrollo (38 pines) & 1 & 8{,}00  & 8{,}00  \\
  & Módulo de relés USB 8 canales 5V DC   & 1 & 12{,}00 & 12{,}00 \\
\midrule
\multirow{3}{*}{Medición}
  & PZEM-004T v3.0 (con TC de pinza)      & 3 & 12{,}00 & 36{,}00 \\
  & \textit{Placas PCB de interfaz (×2)}  & 2 & —       & \textit{a confirmar} \\
  & Cables de medición y puntas           & — & 10{,}00 & 10{,}00 \\
\midrule
\multirow{2}{*}{Potencia}
  & Contactor 12A / 220V AC (2NA+2NC)     & 4 & 9{,}00  & 36{,}00 \\
  & Transformador 220V / 1000--2500V 30mA & 1 & 80{,}00 & 80{,}00 \\
  & Transformador 220V / 2{,}2V / 75VA    & 1 & 35{,}00 & 35{,}00 \\
\midrule
\multirow{3}{*}{Seguridad}
  & Fusible cerámico 10A + portafusible    & 4 & 2{,}00  & 8{,}00  \\
  & Piloto indicador LED 220V (rojo/verde) & 4 & 1{,}50  & 6{,}00  \\
  & Pulsador de paro de emergencia (hongo) & 1 & 8{,}00  & 8{,}00  \\
\midrule
\multirow{2}{*}{Conexionado}
  & Cable H07V-K 1,5~mm² (carrete 10~m)   & 2 & 6{,}00  & 12{,}00 \\
  & Borneras, terminales y bridas          & — & 8{,}00  & 8{,}00  \\
  & Toma corriente industrial + fichas     & 2 & 5{,}00  & 10{,}00 \\
\midrule
Estructura
  & Anvil Flight Case (aluminio)           & 1 & 60{,}00 & 60{,}00 \\
  & Riel DIN + panel frontal aluminio      & 1 & 15{,}00 & 15{,}00 \\
\midrule
\multicolumn{4}{r}{\textbf{Total estimado}} & \textbf{344{,}00} \\
\multicolumn{4}{r}{Margen de contingencia (15\%)} & 51{,}60 \\
\multicolumn{4}{r}{\textbf{Total con contingencia}} & \textbf{395{,}60} \\
\bottomrule
\end{tabular}
\end{table}

\noindent\textit{Nota: el costo de horas de ingeniería y desarrollo de
software no se incluye en el presupuesto del prototipo, dado que
corresponde al esfuerzo académico del proyecto de grado.}

% ------------------------------------------------------------
\subsection{Comparación con soluciones comerciales}
\label{subsec:comparacion_comercial}
% ------------------------------------------------------------

A modo de referencia, la siguiente tabla compara el equipo desarrollado
con equipos comerciales equivalentes disponibles en el mercado nacional,
en particular los modelos de la línea SYGlab (SYG~Group, Rosario, Argentina),
que realizan ensayos similares bajo norma IEC~60335.

\begin{table}[H]
\centering
\caption{Comparación entre el prototipo desarrollado y equipos comerciales equivalentes}
\label{tab:comparacion_comercial}
\renewcommand{\arraystretch}{1.3}
\begin{tabular}{@{} l c c c @{}}
\toprule
\textbf{Característica}       & \textbf{Prototipo (este trabajo)} & \textbf{SYGlab CL12-S0} & \textbf{SYGlab MG5} \\
\midrule
Ensayo de consumo / potencia  & \checkmark & \checkmark & \checkmark \\
Corriente de fuga pasiva      & Arq. prevista & \checkmark & \checkmark \\
Corriente de fuga activa      & Arq. prevista & \checkmark & \checkmark \\
Puesta a tierra               & Arq. prevista & \checkmark & \checkmark \\
Rigidez dieléctrica           & Arq. prevista & \checkmark & \checkmark \\
Clase I y II                  & \checkmark & \checkmark & \checkmark \\
Canales simultáneos           & 1          & 1          & 5          \\
Pantalla / GUI                & PC (PyQt)  & Pantalla HD integrada & Pantalla HD integrada \\
Exportación de datos          & CSV / JSON & Etiqueta QR + base de datos & Etiqueta QR + base de datos \\
Norma declarada               & IEC~60335  & IEC / IRAM & IEC / IRAM \\
Costo estimado                & $\sim$\$400 USD & Consultar  & Consultar  \\
\bottomrule
\end{tabular}
\end{table}

La comparación pone en evidencia que el prototipo desarrollado cubre
el mismo conjunto de ensayos que los equipos comerciales de referencia,
con la diferencia de que en este trabajo se implementa íntegramente
solo el ensayo de consumo, reservando los restantes para etapas
posteriores del proyecto.
La ventaja diferencial del prototipo radica en su bajo costo de
construcción, la apertura total del código fuente y la posibilidad
de adaptación a distintas normativas sin restricciones de licencia.

%
% ============================================================

% ------------------------------------------------------------
\section{Requisitos del Sistema a Desarrollar — Hardware}
\label{sec:requisitos_hw}
% ------------------------------------------------------------

El subsistema de hardware constituye el núcleo físico del equipo de ensayos.
Su función es ejecutar, de forma controlada y repetible, los procedimientos
de medición definidos por la norma IEC~60335, garantizando la seguridad
del operario y la trazabilidad de los resultados.

El dispositivo debe cubrir los siguientes ensayos:

\begin{itemize}
    \item Medición de consumo de potencia (tensión, corriente, potencia activa y factor de potencia).
    \item Medición de corriente de fuga pasiva.
    \item Medición de corriente de fuga activa.
    \item Medición de resistencia de puesta a tierra.
    \item Ensayo de rigidez dieléctrica (tensión resistida).
\end{itemize}

En el marco del presente proyecto se desarrolla el prototipo funcional
centrado en el \textbf{ensayo de consumo de potencia}, siendo los ensayos
restantes contemplados en la arquitectura del sistema para su incorporación
en etapas futuras.

Las medidas de seguridad implementadas en el hardware incluyen:
aislamiento galvánico entre etapas de alta y baja tensión,
desconexión automática ante valores fuera de rango,
indicadores visuales de estado y protocolos de paro de emergencia.

% ------------------------------------------------------------
\section{Metodología y Desarrollo del Hardware}
\label{sec:metodologia_hw}
% ------------------------------------------------------------

\subsection{Diseño del dispositivo}
\label{subsec:diseno_dispositivo}

\subsubsection{Criterios de selección de componentes}
\label{subsubsec:criterios_seleccion}

La selección de los componentes del equipo se realizó mediante un
\textbf{análisis de mercado} en el que se evaluaron los siguientes criterios:

\begin{enumerate}
    \item \textbf{Relación calidad-precio:} se priorizaron componentes que
          ofrecen prestaciones adecuadas para los rangos de medición exigidos
          por la norma IEC~60335, sin incurrir en costos de equipamiento
          industrial de alta gama.

    \item \textbf{Soporte de la comunidad y documentación:} se seleccionaron
          plataformas y módulos con amplia documentación técnica, bibliotecas
          de software activas y comunidades de desarrollo consolidadas, lo que
          reduce el riesgo de obsolescencia y facilita la resolución de
          problemas durante el desarrollo.

    \item \textbf{Disponibilidad en el mercado local e importación:}
          se verificó la disponibilidad de cada componente tanto en
          distribuidores nacionales como en plataformas de importación
          (Mercado Libre, AliExpress, DigiKey), asegurando la reposición
          del stock ante fallas o escalado del proyecto.
\end{enumerate}

\subsubsection{Descripción de componentes seleccionados}
\label{subsubsec:componentes}

A continuación se describen los componentes principales del dispositivo,
agrupados por función.

% - - - - - - - - - - - - - - - - - - - - -
\paragraph{Unidad de procesamiento — ESP32}

El microcontrolador seleccionado es el \textbf{ESP32} (Espressif Systems),
en su versión de módulo de desarrollo de 38 pines.
Se trata de un sistema en chip (SoC) de doble núcleo Xtensa LX6
a 240~MHz, con 520~KB de SRAM, conectividad Wi-Fi 802.11~b/g/n
y Bluetooth 4.2/BLE, y soporte nativo para los protocolos
UART, SPI, I\textsuperscript{2}C y PWM.

Su elección se justifica por la amplia disponibilidad en el mercado
argentino, el extenso ecosistema de bibliotecas compatibles con
Arduino y MicroPython, y por su capacidad de comunicarse simultáneamente
con múltiples módulos periféricos a través de UART, requisito
indispensable para la integración con los módulos PZEM-004T.
Adicionalmente, la conectividad inalámbrica incorporada abre la posibilidad
de comunicación remota en versiones futuras del equipo, sin requerir
hardware adicional.

% - - - - - - - - - - - - - - - - - - - - -
\paragraph{Módulos de medición de potencia — PZEM-004T v3.0}

Para la adquisición de parámetros eléctricos se utilizan tres módulos
\textbf{PZEM-004T v3.0} (Peacefair).
Cada módulo mide en tiempo real:

\begin{itemize}
    \item Tensión de línea: rango 80--260~V\textsubscript{AC},
          resolución 0{,}1~V, precisión $\pm$0{,}5\%.
    \item Corriente: rango 0--100~A (con transformador de corriente externo
          de pinza incluido), resolución 0{,}001~A, precisión $\pm$0{,}5\%.
    \item Potencia activa: rango 0--23~kW, resolución 0{,}1~W.
    \item Energía acumulada: rango 0--9999~kWh, resolución 1~Wh.
    \item Frecuencia: rango 45--65~Hz, resolución 0{,}1~Hz.
    \item Factor de potencia: rango 0{,}00--1{,}00, resolución 0{,}01.
\end{itemize}

La comunicación con el ESP32 se realiza mediante el protocolo
\textbf{UART/Modbus RTU} a 9600~baudios.
Cada módulo posee una dirección Modbus configurable (0x01--0xF7),
lo que permite conectar los tres módulos en un bus RS-485 compartido
y direccionarlos de forma individual desde el microcontrolador.

La elección de tres unidades responde a la arquitectura de conmutación
del equipo: uno para el ensayo de consumo nominal,
uno para la medición de fuga pasiva
y uno para la medición de fuga activa,
permitiendo la transición entre ensayos sin necesidad de reconexión manual.

% - - - - - - - - - - - - - - - - - - - - -
\paragraph{Sistema de conmutación — Contactores y módulo de relés}

La conmutación entre ensayos se realiza mediante dos tipos de actuadores:

\textbf{Contactores monofásicos 12~A (×4):}
se emplean contactores de 12~A con bobina de 220~V~AC,
configuración de contactos 2~NA + 2~NC.
Su función es la conmutación de la tensión de red (220~V~AC)
hacia las distintas etapas del equipo (alimentación del DUT,
circuito de rigidez dieléctrica, circuito de fuga).
La acción mecánica de apertura garantiza el aislamiento galvánico
efectivo entre etapas, lo que es un requisito de seguridad
fundamental cuando se trabaja con tensiones de ensayo elevadas
(hasta 2500~V en el ensayo de rigidez dieléctrica).

\textbf{Módulo de relés USB de 8 canales (×1):}
se utiliza un módulo de relés de 8 canales controlado por USB,
con relés de 5~V~DC y capacidad de conmutación de hasta 10~A~AC.
Este módulo gestiona señales de control de baja potencia:
activación de indicadores, habilitación de etapas auxiliares
y señalización de estado al software.
La comunicación por USB permite el control directo desde la PC
sin necesidad de pasar por el ESP32, simplificando la arquitectura
de control para funciones no críticas en tiempo.

% - - - - - - - - - - - - - - - - - - - - -
\paragraph{Transformadores de ensayo}

El equipo incorpora dos transformadores de propósito específico:

\textbf{Transformador de rigidez dieléctrica — 220~V / 1000--2500~V:}
transformador elevador con salida seleccionable entre 1000~V
(ensayo Clase~I) y 2500~V (ensayo Clase~II),
con corriente de cortocircuito limitada a 30~mA.
Esta limitación es un requisito de seguridad estipulado en la
norma IEC~60335: la corriente máxima admisible durante el ensayo
de tensión resistida es de 50~mA, y la restricción a 30~mA en
la fuente asegura que una falla del aislamiento no derive en
una corriente destructiva para el DUT ni peligrosa para el operario.

\textbf{Transformador de puesta a tierra — 220~V / 2{,}2~V / 75~VA:}
transformador reductor de 75~VA utilizado en el ensayo de
continuidad de tierra.
Entrega una tensión de 2{,}2~V~AC con capacidad de corriente
de ensayo de hasta 15~A~AC (equivalente a $75~\text{VA} / 5~\text{V} \approx 15~\text{A}$),
cumpliendo el requisito de la norma de inyectar al menos 25~A
o el 200\% de la corriente nominal del fusible de protección,
según el método aplicado.
La baja tensión de salida garantiza la seguridad del operario
durante el ensayo.

% - - - - - - - - - - - - - - - - - - - - -
\paragraph{Carcasa y estructura — Anvil Flight Case}

El equipo se aloja en un \textbf{Anvil Flight Case} de aluminio y
madera contrachapada con herrajes de acero.
Este tipo de contenedor, habitualmente utilizado en equipamiento
de audio y escenografía profesional, presenta las siguientes ventajas
para la aplicación:

\begin{itemize}
    \item Rigidez estructural adecuada para soportar transformadores
          y contactores sin deformaciones.
    \item Cierre con pestillos de seguridad que evitan apertura accidental.
    \item Carácter portátil con manija y esquinas reforzadas,
          compatible con el concepto de equipo de laboratorio portátil.
    \item Posibilidad de montaje de paneles frontales y rieles DIN
          en el interior para organizar los componentes.
\end{itemize}

% - - - - - - - - - - - - - - - - - - - - -
\paragraph{Placas de interfaz y conexionado}

\noindent\textbf{Nota para revisión:} \textit{Se incorporarán dos placas de
circuito impreso (PCB) de diseño propio cuyas especificaciones
se encuentran en proceso de definición (a confirmar con Nico).
Se completará esta sección con: función de cada placa,
dimensiones, descripción del circuito y justificación de diseño.}

% ------------------------------------------------------------
\subsection{Arquitectura del dispositivo}
\label{subsec:arquitectura_hw}
% ------------------------------------------------------------

La arquitectura del dispositivo se organiza en cinco bloques funcionales
interconectados a través de un bus de comunicación principal
gestionado por el ESP32:

\begin{enumerate}
    \item \textbf{Bloque de potencia:} transformadores, contactores y
          protecciones. Maneja las tensiones de red y de ensayo.
    \item \textbf{Bloque de adquisición:} módulos PZEM-004T. Realizan
          la medición de los parámetros eléctricos del DUT.
    \item \textbf{Bloque de control:} ESP32 y módulo de relés USB.
          Orchestran la secuencia de ensayos y la comunicación.
    \item \textbf{Bloque de comunicación:} interfaz USB/UART hacia la PC.
          Transmite datos crudos y recibe comandos de la interfaz gráfica.
    \item \textbf{Bloque de seguridad:} fusibles, protecciones térmicas
          e indicadores. Actúan de forma independiente ante condiciones
          de falla.
\end{enumerate}

% (Figura del diagrama de bloques — reemplazar con tu figura real)
\begin{figure}[H]
    \centering
    % \includegraphics[width=0.85\textwidth]{figuras/diagrama_bloques_hw.png}
    \fbox{\parbox{0.85\textwidth}{\centering\vspace{2cm}
    \textit{[Insertar aquí el diagrama de bloques del hardware]}
    \vspace{2cm}}}
    \caption{Diagrama de bloques del dispositivo de ensayos.}
    \label{fig:diagrama_bloques_hw}
\end{figure}

% ------------------------------------------------------------
\subsection{Protocolo de comunicación}
\label{subsec:protocolo_comunicacion}
% ------------------------------------------------------------

La comunicación entre componentes del equipo sigue el esquema
descrito a continuación:

\begin{itemize}
    \item \textbf{ESP32 $\leftrightarrow$ módulos PZEM-004T:}
          protocolo Modbus~RTU sobre UART a 9600~baudios, 8N1
          (8 bits de datos, sin paridad, 1 bit de parada).
          Los tres módulos comparten el bus serie;
          cada uno responde únicamente a su dirección Modbus asignada
          (0x01, 0x02 y 0x03 respectivamente).

    \item \textbf{ESP32 $\leftrightarrow$ PC:}
          puerto UART expuesto como dispositivo CDC-ACM (serial virtual)
          a través del chip USB-to-UART integrado en el módulo de desarrollo
          (CH340 o CP2102 según la revisión del módulo).
          La velocidad de comunicación es de 115200~baudios.
          El protocolo de capa de aplicación es un formato JSON
          delimitado por salto de línea (\texttt{newline-delimited JSON}),
          elegido por su legibilidad, facilidad de parseo en Python
          y extensibilidad para agregar nuevos tipos de medición.

    \item \textbf{PC $\leftrightarrow$ módulo de relés USB:}
          protocolo propietario del módulo sobre USB HID / CDC,
          controlado directamente desde el software en Python
          mediante la biblioteca \texttt{pyserial} o el driver
          correspondiente al chip del módulo.
\end{itemize}

% ------------------------------------------------------------
\section{Evaluación de costos}
\label{sec:costos}
% ------------------------------------------------------------

La siguiente tabla presenta el desglose de costos del prototipo
desarrollado. Los precios expresados en dólares estadounidenses (USD)
corresponden a valores de referencia de mercado obtenidos en
distribuidores nacionales y plataformas de importación
(Mercado Libre Argentina, AliExpress, DigiKey),
relevados durante el segundo semestre de 2025.
Los valores deben considerarse estimativos;
variaciones de tipo de cambio o disponibilidad de stock
pueden afectar el costo final.

\begin{table}[H]
\centering
\caption{Desglose de costos del prototipo — valores de referencia en USD}
\label{tab:costos_detalle}
\renewcommand{\arraystretch}{1.3}
\begin{tabular}{@{} l l r r r @{}}
\toprule
\textbf{Categoría} & \textbf{Componente} & \textbf{Cant.} & \textbf{P. Unit. (USD)} & \textbf{Subtotal (USD)} \\
\midrule
\multirow{2}{*}{Procesamiento}
  & ESP32 módulo de desarrollo (38 pines) & 1 & 8{,}00  & 8{,}00  \\
  & Módulo de relés USB 8 canales 5V DC   & 1 & 12{,}00 & 12{,}00 \\
\midrule
\multirow{3}{*}{Medición}
  & PZEM-004T v3.0 (con TC de pinza)      & 3 & 12{,}00 & 36{,}00 \\
  & \textit{Placas PCB de interfaz (×2)}  & 2 & —       & \textit{a confirmar} \\
  & Cables de medición y puntas           & — & 10{,}00 & 10{,}00 \\
\midrule
\multirow{2}{*}{Potencia}
  & Contactor 12A / 220V AC (2NA+2NC)     & 4 & 9{,}00  & 36{,}00 \\
  & Transformador 220V / 1000--2500V 30mA & 1 & 80{,}00 & 80{,}00 \\
  & Transformador 220V / 2{,}2V / 75VA    & 1 & 35{,}00 & 35{,}00 \\
\midrule
\multirow{3}{*}{Seguridad}
  & Fusible cerámico 10A + portafusible    & 4 & 2{,}00  & 8{,}00  \\
  & Piloto indicador LED 220V (rojo/verde) & 4 & 1{,}50  & 6{,}00  \\
  & Pulsador de paro de emergencia (hongo) & 1 & 8{,}00  & 8{,}00  \\
\midrule
\multirow{2}{*}{Conexionado}
  & Cable H07V-K 1,5~mm² (carrete 10~m)   & 2 & 6{,}00  & 12{,}00 \\
  & Borneras, terminales y bridas          & — & 8{,}00  & 8{,}00  \\
  & Toma corriente industrial + fichas     & 2 & 5{,}00  & 10{,}00 \\
\midrule
Estructura
  & Anvil Flight Case (aluminio)           & 1 & 60{,}00 & 60{,}00 \\
  & Riel DIN + panel frontal aluminio      & 1 & 15{,}00 & 15{,}00 \\
\midrule
\multicolumn{4}{r}{\textbf{Total estimado}} & \textbf{344{,}00} \\
\multicolumn{4}{r}{Margen de contingencia (15\%)} & 51{,}60 \\
\multicolumn{4}{r}{\textbf{Total con contingencia}} & \textbf{395{,}60} \\
\bottomrule
\end{tabular}
\end{table}

\noindent\textit{Nota: el costo de horas de ingeniería y desarrollo de
software no se incluye en el presupuesto del prototipo, dado que
corresponde al esfuerzo académico del proyecto de grado.}

% ------------------------------------------------------------
\subsection{Comparación con soluciones comerciales}
\label{subsec:comparacion_comercial}
% ------------------------------------------------------------

A modo de referencia, la siguiente tabla compara el equipo desarrollado
con equipos comerciales equivalentes disponibles en el mercado nacional,
en particular los modelos de la línea SYGlab (SYG~Group, Rosario, Argentina),
que realizan ensayos similares bajo norma IEC~60335.

\begin{table}[H]
\centering
\caption{Comparación entre el prototipo desarrollado y equipos comerciales equivalentes}
\label{tab:comparacion_comercial}
\renewcommand{\arraystretch}{1.3}
\begin{tabular}{@{} l c c c @{}}
\toprule
\textbf{Característica}       & \textbf{Prototipo (este trabajo)} & \textbf{SYGlab CL12-S0} & \textbf{SYGlab MG5} \\
\midrule
Ensayo de consumo / potencia  & \checkmark & \checkmark & \checkmark \\
Corriente de fuga pasiva      & Arq. prevista & \checkmark & \checkmark \\
Corriente de fuga activa      & Arq. prevista & \checkmark & \checkmark \\
Puesta a tierra               & Arq. prevista & \checkmark & \checkmark \\
Rigidez dieléctrica           & Arq. prevista & \checkmark & \checkmark \\
Clase I y II                  & \checkmark & \checkmark & \checkmark \\
Canales simultáneos           & 1          & 1          & 5          \\
Pantalla / GUI                & PC (PyQt)  & Pantalla HD integrada & Pantalla HD integrada \\
Exportación de datos          & CSV / JSON & Etiqueta QR + base de datos & Etiqueta QR + base de datos \\
Norma declarada               & IEC~60335  & IEC / IRAM & IEC / IRAM \\
Costo estimado                & $\sim$\$400 USD & Consultar  & Consultar  \\
\bottomrule
\end{tabular}
\end{table}

La comparación pone en evidencia que el prototipo desarrollado cubre
el mismo conjunto de ensayos que los equipos comerciales de referencia,
con la diferencia de que en este trabajo se implementa íntegramente
solo el ensayo de consumo, reservando los restantes para etapas
posteriores del proyecto.
La ventaja diferencial del prototipo radica en su bajo costo de
construcción, la apertura total del código fuente y la posibilidad
de adaptación a distintas normativas sin restricciones de licencia.

%
% ============================================================

% ------------------------------------------------------------
\section{Requisitos del Sistema a Desarrollar — Hardware}
\label{sec:requisitos_hw}
% ------------------------------------------------------------

El subsistema de hardware constituye el núcleo físico del equipo de ensayos.
Su función es ejecutar, de forma controlada y repetible, los procedimientos
de medición definidos por la norma IEC~60335, garantizando la seguridad
del operario y la trazabilidad de los resultados.

El dispositivo debe cubrir los siguientes ensayos:

\begin{itemize}
    \item Medición de consumo de potencia (tensión, corriente, potencia activa y factor de potencia).
    \item Medición de corriente de fuga pasiva.
    \item Medición de corriente de fuga activa.
    \item Medición de resistencia de puesta a tierra.
    \item Ensayo de rigidez dieléctrica (tensión resistida).
\end{itemize}

En el marco del presente proyecto se desarrolla el prototipo funcional
centrado en el \textbf{ensayo de consumo de potencia}, siendo los ensayos
restantes contemplados en la arquitectura del sistema para su incorporación
en etapas futuras.

Las medidas de seguridad implementadas en el hardware incluyen:
aislamiento galvánico entre etapas de alta y baja tensión,
desconexión automática ante valores fuera de rango,
indicadores visuales de estado y protocolos de paro de emergencia.

% ------------------------------------------------------------
\section{Metodología y Desarrollo del Hardware}
\label{sec:metodologia_hw}
% ------------------------------------------------------------

\subsection{Diseño del dispositivo}
\label{subsec:diseno_dispositivo}

\subsubsection{Criterios de selección de componentes}
\label{subsubsec:criterios_seleccion}

La selección de los componentes del equipo se realizó mediante un
\textbf{análisis de mercado} en el que se evaluaron los siguientes criterios:

\begin{enumerate}
    \item \textbf{Relación calidad-precio:} se priorizaron componentes que
          ofrecen prestaciones adecuadas para los rangos de medición exigidos
          por la norma IEC~60335, sin incurrir en costos de equipamiento
          industrial de alta gama.

    \item \textbf{Soporte de la comunidad y documentación:} se seleccionaron
          plataformas y módulos con amplia documentación técnica, bibliotecas
          de software activas y comunidades de desarrollo consolidadas, lo que
          reduce el riesgo de obsolescencia y facilita la resolución de
          problemas durante el desarrollo.

    \item \textbf{Disponibilidad en el mercado local e importación:}
          se verificó la disponibilidad de cada componente tanto en
          distribuidores nacionales como en plataformas de importación
          (Mercado Libre, AliExpress, DigiKey), asegurando la reposición
          del stock ante fallas o escalado del proyecto.
\end{enumerate}

\subsubsection{Descripción de componentes seleccionados}
\label{subsubsec:componentes}

A continuación se describen los componentes principales del dispositivo,
agrupados por función.

% - - - - - - - - - - - - - - - - - - - - -
\paragraph{Unidad de procesamiento — ESP32}

El microcontrolador seleccionado es el \textbf{ESP32} (Espressif Systems),
en su versión de módulo de desarrollo de 38 pines.
Se trata de un sistema en chip (SoC) de doble núcleo Xtensa LX6
a 240~MHz, con 520~KB de SRAM, conectividad Wi-Fi 802.11~b/g/n
y Bluetooth 4.2/BLE, y soporte nativo para los protocolos
UART, SPI, I\textsuperscript{2}C y PWM.

Su elección se justifica por la amplia disponibilidad en el mercado
argentino, el extenso ecosistema de bibliotecas compatibles con
Arduino y MicroPython, y por su capacidad de comunicarse simultáneamente
con múltiples módulos periféricos a través de UART, requisito
indispensable para la integración con los módulos PZEM-004T.
Adicionalmente, la conectividad inalámbrica incorporada abre la posibilidad
de comunicación remota en versiones futuras del equipo, sin requerir
hardware adicional.

% - - - - - - - - - - - - - - - - - - - - -
\paragraph{Módulos de medición de potencia — PZEM-004T v3.0}

Para la adquisición de parámetros eléctricos se utilizan tres módulos
\textbf{PZEM-004T v3.0} (Peacefair).
Cada módulo mide en tiempo real:

\begin{itemize}
    \item Tensión de línea: rango 80--260~V\textsubscript{AC},
          resolución 0{,}1~V, precisión $\pm$0{,}5\%.
    \item Corriente: rango 0--100~A (con transformador de corriente externo
          de pinza incluido), resolución 0{,}001~A, precisión $\pm$0{,}5\%.
    \item Potencia activa: rango 0--23~kW, resolución 0{,}1~W.
    \item Energía acumulada: rango 0--9999~kWh, resolución 1~Wh.
    \item Frecuencia: rango 45--65~Hz, resolución 0{,}1~Hz.
    \item Factor de potencia: rango 0{,}00--1{,}00, resolución 0{,}01.
\end{itemize}

La comunicación con el ESP32 se realiza mediante el protocolo
\textbf{UART/Modbus RTU} a 9600~baudios.
Cada módulo posee una dirección Modbus configurable (0x01--0xF7),
lo que permite conectar los tres módulos en un bus RS-485 compartido
y direccionarlos de forma individual desde el microcontrolador.

La elección de tres unidades responde a la arquitectura de conmutación
del equipo: uno para el ensayo de consumo nominal,
uno para la medición de fuga pasiva
y uno para la medición de fuga activa,
permitiendo la transición entre ensayos sin necesidad de reconexión manual.

% - - - - - - - - - - - - - - - - - - - - -
\paragraph{Sistema de conmutación — Contactores y módulo de relés}

La conmutación entre ensayos se realiza mediante dos tipos de actuadores:

\textbf{Contactores monofásicos 12~A (×4):}
se emplean contactores de 12~A con bobina de 220~V~AC,
configuración de contactos 2~NA + 2~NC.
Su función es la conmutación de la tensión de red (220~V~AC)
hacia las distintas etapas del equipo (alimentación del DUT,
circuito de rigidez dieléctrica, circuito de fuga).
La acción mecánica de apertura garantiza el aislamiento galvánico
efectivo entre etapas, lo que es un requisito de seguridad
fundamental cuando se trabaja con tensiones de ensayo elevadas
(hasta 2500~V en el ensayo de rigidez dieléctrica).

\textbf{Módulo de relés USB de 8 canales (×1):}
se utiliza un módulo de relés de 8 canales controlado por USB,
con relés de 5~V~DC y capacidad de conmutación de hasta 10~A~AC.
Este módulo gestiona señales de control de baja potencia:
activación de indicadores, habilitación de etapas auxiliares
y señalización de estado al software.
La comunicación por USB permite el control directo desde la PC
sin necesidad de pasar por el ESP32, simplificando la arquitectura
de control para funciones no críticas en tiempo.

% - - - - - - - - - - - - - - - - - - - - -
\paragraph{Transformadores de ensayo}

El equipo incorpora dos transformadores de propósito específico:

\textbf{Transformador de rigidez dieléctrica — 220~V / 1000--2500~V:}
transformador elevador con salida seleccionable entre 1000~V
(ensayo Clase~I) y 2500~V (ensayo Clase~II),
con corriente de cortocircuito limitada a 30~mA.
Esta limitación es un requisito de seguridad estipulado en la
norma IEC~60335: la corriente máxima admisible durante el ensayo
de tensión resistida es de 50~mA, y la restricción a 30~mA en
la fuente asegura que una falla del aislamiento no derive en
una corriente destructiva para el DUT ni peligrosa para el operario.

\textbf{Transformador de puesta a tierra — 220~V / 2{,}2~V / 75~VA:}
transformador reductor de 75~VA utilizado en el ensayo de
continuidad de tierra.
Entrega una tensión de 2{,}2~V~AC con capacidad de corriente
de ensayo de hasta 15~A~AC (equivalente a $75~\text{VA} / 5~\text{V} \approx 15~\text{A}$),
cumpliendo el requisito de la norma de inyectar al menos 25~A
o el 200\% de la corriente nominal del fusible de protección,
según el método aplicado.
La baja tensión de salida garantiza la seguridad del operario
durante el ensayo.

% - - - - - - - - - - - - - - - - - - - - -
\paragraph{Carcasa y estructura — Anvil Flight Case}

El equipo se aloja en un \textbf{Anvil Flight Case} de aluminio y
madera contrachapada con herrajes de acero.
Este tipo de contenedor, habitualmente utilizado en equipamiento
de audio y escenografía profesional, presenta las siguientes ventajas
para la aplicación:

\begin{itemize}
    \item Rigidez estructural adecuada para soportar transformadores
          y contactores sin deformaciones.
    \item Cierre con pestillos de seguridad que evitan apertura accidental.
    \item Carácter portátil con manija y esquinas reforzadas,
          compatible con el concepto de equipo de laboratorio portátil.
    \item Posibilidad de montaje de paneles frontales y rieles DIN
          en el interior para organizar los componentes.
\end{itemize}

% - - - - - - - - - - - - - - - - - - - - -
\paragraph{Placas de interfaz y conexionado}

\noindent\textbf{Nota para revisión:} \textit{Se incorporarán dos placas de
circuito impreso (PCB) de diseño propio cuyas especificaciones
se encuentran en proceso de definición (a confirmar con Nico).
Se completará esta sección con: función de cada placa,
dimensiones, descripción del circuito y justificación de diseño.}

% ------------------------------------------------------------
\subsection{Arquitectura del dispositivo}
\label{subsec:arquitectura_hw}
% ------------------------------------------------------------

La arquitectura del dispositivo se organiza en cinco bloques funcionales
interconectados a través de un bus de comunicación principal
gestionado por el ESP32:

\begin{enumerate}
    \item \textbf{Bloque de potencia:} transformadores, contactores y
          protecciones. Maneja las tensiones de red y de ensayo.
    \item \textbf{Bloque de adquisición:} módulos PZEM-004T. Realizan
          la medición de los parámetros eléctricos del DUT.
    \item \textbf{Bloque de control:} ESP32 y módulo de relés USB.
          Orchestran la secuencia de ensayos y la comunicación.
    \item \textbf{Bloque de comunicación:} interfaz USB/UART hacia la PC.
          Transmite datos crudos y recibe comandos de la interfaz gráfica.
    \item \textbf{Bloque de seguridad:} fusibles, protecciones térmicas
          e indicadores. Actúan de forma independiente ante condiciones
          de falla.
\end{enumerate}

% (Figura del diagrama de bloques — reemplazar con tu figura real)
\begin{figure}[H]
    \centering
    % \includegraphics[width=0.85\textwidth]{figuras/diagrama_bloques_hw.png}
    \fbox{\parbox{0.85\textwidth}{\centering\vspace{2cm}
    \textit{[Insertar aquí el diagrama de bloques del hardware]}
    \vspace{2cm}}}
    \caption{Diagrama de bloques del dispositivo de ensayos.}
    \label{fig:diagrama_bloques_hw}
\end{figure}

% ------------------------------------------------------------
\subsection{Protocolo de comunicación}
\label{subsec:protocolo_comunicacion}
% ------------------------------------------------------------

La comunicación entre componentes del equipo sigue el esquema
descrito a continuación:

\begin{itemize}
    \item \textbf{ESP32 $\leftrightarrow$ módulos PZEM-004T:}
          protocolo Modbus~RTU sobre UART a 9600~baudios, 8N1
          (8 bits de datos, sin paridad, 1 bit de parada).
          Los tres módulos comparten el bus serie;
          cada uno responde únicamente a su dirección Modbus asignada
          (0x01, 0x02 y 0x03 respectivamente).

    \item \textbf{ESP32 $\leftrightarrow$ PC:}
          puerto UART expuesto como dispositivo CDC-ACM (serial virtual)
          a través del chip USB-to-UART integrado en el módulo de desarrollo
          (CH340 o CP2102 según la revisión del módulo).
          La velocidad de comunicación es de 115200~baudios.
          El protocolo de capa de aplicación es un formato JSON
          delimitado por salto de línea (\texttt{newline-delimited JSON}),
          elegido por su legibilidad, facilidad de parseo en Python
          y extensibilidad para agregar nuevos tipos de medición.

    \item \textbf{PC $\leftrightarrow$ módulo de relés USB:}
          protocolo propietario del módulo sobre USB HID / CDC,
          controlado directamente desde el software en Python
          mediante la biblioteca \texttt{pyserial} o el driver
          correspondiente al chip del módulo.
\end{itemize}

% ------------------------------------------------------------
\section{Evaluación de costos}
\label{sec:costos}
% ------------------------------------------------------------

La siguiente tabla presenta el desglose de costos del prototipo
desarrollado. Los precios expresados en dólares estadounidenses (USD)
corresponden a valores de referencia de mercado obtenidos en
distribuidores nacionales y plataformas de importación
(Mercado Libre Argentina, AliExpress, DigiKey),
relevados durante el segundo semestre de 2025.
Los valores deben considerarse estimativos;
variaciones de tipo de cambio o disponibilidad de stock
pueden afectar el costo final.

\begin{table}[H]
\centering
\caption{Desglose de costos del prototipo — valores de referencia en USD}
\label{tab:costos_detalle}
\renewcommand{\arraystretch}{1.3}
\begin{tabular}{@{} l l r r r @{}}
\toprule
\textbf{Categoría} & \textbf{Componente} & \textbf{Cant.} & \textbf{P. Unit. (USD)} & \textbf{Subtotal (USD)} \\
\midrule
\multirow{2}{*}{Procesamiento}
  & ESP32 módulo de desarrollo (38 pines) & 1 & 8{,}00  & 8{,}00  \\
  & Módulo de relés USB 8 canales 5V DC   & 1 & 12{,}00 & 12{,}00 \\
\midrule
\multirow{3}{*}{Medición}
  & PZEM-004T v3.0 (con TC de pinza)      & 3 & 12{,}00 & 36{,}00 \\
  & \textit{Placas PCB de interfaz (×2)}  & 2 & —       & \textit{a confirmar} \\
  & Cables de medición y puntas           & — & 10{,}00 & 10{,}00 \\
\midrule
\multirow{2}{*}{Potencia}
  & Contactor 12A / 220V AC (2NA+2NC)     & 4 & 9{,}00  & 36{,}00 \\
  & Transformador 220V / 1000--2500V 30mA & 1 & 80{,}00 & 80{,}00 \\
  & Transformador 220V / 2{,}2V / 75VA    & 1 & 35{,}00 & 35{,}00 \\
\midrule
\multirow{3}{*}{Seguridad}
  & Fusible cerámico 10A + portafusible    & 4 & 2{,}00  & 8{,}00  \\
  & Piloto indicador LED 220V (rojo/verde) & 4 & 1{,}50  & 6{,}00  \\
  & Pulsador de paro de emergencia (hongo) & 1 & 8{,}00  & 8{,}00  \\
\midrule
\multirow{2}{*}{Conexionado}
  & Cable H07V-K 1,5~mm² (carrete 10~m)   & 2 & 6{,}00  & 12{,}00 \\
  & Borneras, terminales y bridas          & — & 8{,}00  & 8{,}00  \\
  & Toma corriente industrial + fichas     & 2 & 5{,}00  & 10{,}00 \\
\midrule
Estructura
  & Anvil Flight Case (aluminio)           & 1 & 60{,}00 & 60{,}00 \\
  & Riel DIN + panel frontal aluminio      & 1 & 15{,}00 & 15{,}00 \\
\midrule
\multicolumn{4}{r}{\textbf{Total estimado}} & \textbf{344{,}00} \\
\multicolumn{4}{r}{Margen de contingencia (15\%)} & 51{,}60 \\
\multicolumn{4}{r}{\textbf{Total con contingencia}} & \textbf{395{,}60} \\
\bottomrule
\end{tabular}
\end{table}

\noindent\textit{Nota: el costo de horas de ingeniería y desarrollo de
software no se incluye en el presupuesto del prototipo, dado que
corresponde al esfuerzo académico del proyecto de grado.}

% ------------------------------------------------------------
\subsection{Comparación con soluciones comerciales}
\label{subsec:comparacion_comercial}
% ------------------------------------------------------------

A modo de referencia, la siguiente tabla compara el equipo desarrollado
con equipos comerciales equivalentes disponibles en el mercado nacional,
en particular los modelos de la línea SYGlab (SYG~Group, Rosario, Argentina),
que realizan ensayos similares bajo norma IEC~60335.

\begin{table}[H]
\centering
\caption{Comparación entre el prototipo desarrollado y equipos comerciales equivalentes}
\label{tab:comparacion_comercial}
\renewcommand{\arraystretch}{1.3}
\begin{tabular}{@{} l c c c @{}}
\toprule
\textbf{Característica}       & \textbf{Prototipo (este trabajo)} & \textbf{SYGlab CL12-S0} & \textbf{SYGlab MG5} \\
\midrule
Ensayo de consumo / potencia  & \checkmark & \checkmark & \checkmark \\
Corriente de fuga pasiva      & Arq. prevista & \checkmark & \checkmark \\
Corriente de fuga activa      & Arq. prevista & \checkmark & \checkmark \\
Puesta a tierra               & Arq. prevista & \checkmark & \checkmark \\
Rigidez dieléctrica           & Arq. prevista & \checkmark & \checkmark \\
Clase I y II                  & \checkmark & \checkmark & \checkmark \\
Canales simultáneos           & 1          & 1          & 5          \\
Pantalla / GUI                & PC (PyQt)  & Pantalla HD integrada & Pantalla HD integrada \\
Exportación de datos          & CSV / JSON & Etiqueta QR + base de datos & Etiqueta QR + base de datos \\
Norma declarada               & IEC~60335  & IEC / IRAM & IEC / IRAM \\
Costo estimado                & $\sim$\$400 USD & Consultar  & Consultar  \\
\bottomrule
\end{tabular}
\end{table}

La comparación pone en evidencia que el prototipo desarrollado cubre
el mismo conjunto de ensayos que los equipos comerciales de referencia,
con la diferencia de que en este trabajo se implementa íntegramente
solo el ensayo de consumo, reservando los restantes para etapas
posteriores del proyecto.
La ventaja diferencial del prototipo radica en su bajo costo de
construcción, la apertura total del código fuente y la posibilidad
de adaptación a distintas normativas sin restricciones de licencia.
